\documentclass[a4paper,12pt]{article}
\usepackage[spanish]{babel}
\addto\captionsspanish{\renewcommand{\tablename}{Tabla}}
\usepackage[utf8]{inputenc}
\usepackage[T1]{fontenc}
\usepackage{graphicx}
\usepackage{amsmath}
\usepackage{amssymb}
\usepackage{geometry}
\usepackage{enumitem}
\usepackage{xcolor}
\usepackage{hyperref}
\usepackage{multirow}
\usepackage{comment}

% Improved typography and tables
\usepackage{microtype}
\usepackage{booktabs}
\usepackage{caption}
\usepackage{siunitx}
\usepackage{fancyhdr}
\usepackage{float}
\usepackage{circuitikz}
\usetikzlibrary{babel, positioning, decorations.pathmorphing}

% Configuration
\hypersetup{
    colorlinks=true,
    linkcolor=blue!70!black,
    citecolor=blue!70!black,
    urlcolor=blue!70!black
}

\geometry{left=2.5cm, right=2.5cm, top=2.5cm, bottom=2.5cm}

\sisetup{
    output-decimal-marker = {,},
    per-mode = symbol
}

% Header and Footer
\pagestyle{fancy}
\fancyhf{}
\fancyhead[L]{PRACTICA 2 DE DISPS ELECTRóNICOS: ESTUDIO DEL TRANSISTOR}
\fancyhead[R]{\thepage}
\renewcommand{\headrulewidth}{0.4pt}

\begin{document}

% Custom Cover Page
\begin{titlepage}
    \centering
    \vspace*{2cm}

    {\LARGE\bfseries UDO/NúCLEO DE SUCRE \par}
    \vspace{0.5cm}
    {\large Departamento de Física\par}
    {\large Laboratorio de Electrónica\par}

    \vspace{3cm}

    {\Huge\bfseries PRáCTICA 2\par}
    \vspace{1cm}
    {\LARGE\bfseries ESTUDIO DEL TRANSISTOR\par}

    \vfill

    \begin{figure}[h]
        \centering
        \begin{circuitikz}[scale=1.5]
            % PNP Transistor Symbol
            \draw (0,0) node[pnp, scale=2] (Q1) {};
            \node[above=0.3cm] at (Q1.E) {\large E};
            \node[below=0.3cm] at (Q1.C) {\large C};
            \node[left=0.3cm] at (Q1.B) {\large B};
        \end{circuitikz}
    \end{figure}

    \vfill

    {\large Laboratorio de Electrónica\par}
    \vspace{1cm}
    {\large \today\par}

\end{titlepage}

\tableofcontents
\newpage

\listoffigures
\newpage

\listoftables
\newpage


\section{Objetivos}
\begin{enumerate}
    \item Conocer los fundamentos del transistor.
    \item Medir los efectos sobre la corriente emisor-base de la polarización directa (normal) y
          de la polarización inversa en el circuito emisor-base.
    \item Medir los efectos sobre la corriente de colector de las polarizaciones directa e inversa
          en el circuito emisor-base.
    \item Medir \( I_{CBO} \).
\end{enumerate}



\section{Información Preliminar}

\subsection{El transistor: dispositivo de tres elementos}

Hasta ahora hemos tratado de los dispositivos electrónicos de dos elementos o electrodos: los diodos
de semiconductores. Hemos encontrado una diversidad de diodos con sus características
peculiares, que sirven para diferentes propósitos. En nuestros experimentos hemos empleado diodos como
rectificadores de corriente alterna, como reguladores de tensión y corriente, y como limitadores. En esta
Práctica vamos a conocer un dispositivo electrónico de tres elementos, el transistor, cuyas características
amplían considerablemente el campo de aplicaciones de la electrónica.

El interés científico de los semiconductores condujo al invento del transistor. Este dispositivo
semiconductor puede realizar funciones tales como amplificación, detección y oscilación. Una ventaja
del transistor es que es pequeño y ligero, lo que permite la miniaturización del equipo electrónico,
y además es sólido. Por lo tanto, la microfonía no constituye un problema como en los tubos de vacío.
No tiene filamentos y por consiguiente no requiere potencia para filamento. El transistor funciona con
tensiones de alimentación bajas y utiliza poca potencia. No requiere período alguno de calentamiento y
funciona tan pronto como se le aplica la potencia. Un transistor tiene menos conexiones y por consiguiente
sus circuitos son básicamente más sencillos que los de un tubo. Un inconveniente de los transistores es
su sensibilidad al calor, pero en cambio una ventaja es que no generan tanto calor como los tubos.

Existen muchos tipos de transistores que pueden ser clasificados según el material básico con que están
formados. En esta categoría encontramos los transistores de germanio y de silicio. También pueden ser
clasificados los transistores de acuerdo con su proceso de construcción. En esta categoría encontramos
varios tipos de transistores de unión o contacto tales como los fabricados por el proceso llamado de
crecimiento, contacto de aleación, de efecto de campo, mesa, mesa epitaxial y los transistores planares
y de puntas de contacto. Los transistores se pueden clasificar también por el número de sus elementos.
Entonces tenemos el tríodo o transistor de tres elementos y el tetrodo o transistor de cuatro elementos.
Por último se pueden clasificar de acuerdo con su aptitud para disipar potencias. En cuanto a esto existe
un amplio margen, desde el tipo de baja potencia (menos de \SI{50}{\milli\watt}) hasta el de alta potencia
\SI{2}{\watt} y más).

Los transistores se fabrican en diferentes formas y tamaños (fig.~\ref{fig:11-1}). También hay variantes
en cuanto a las disposiciones de la base y la manera de montar los transistores en el circuito.

\begin{figure}[h]
    \centering
    \begin{tikzpicture}[scale=1.2, transform shape]
        % TO-92 Package (Plastic)
        \begin{scope}[shift={(0,0)}]
            % Leads
            \draw[gray, thick] (-0.3, -1.5) -- (-0.3, 0);
            \draw[gray, thick] (0, -1.5) -- (0, 0);
            \draw[gray, thick] (0.3, -1.5) -- (0.3, 0);

            % Body (D-shape cylinder approximation)
            \draw[fill=black!80] (-0.5, 0) arc (180:0:0.5) -- (0.5, 0) -- (0.5, 0.8) arc (0:180:0.5) -- cycle;
            \draw[fill=black!70] (-0.5, 0) -- (0.5, 0) -- (0.5, -0.1) -- (-0.5, -0.1) -- cycle; % Bottom
            \fill[black!90] (-0.5,0) rectangle (0.5, 0.8);
            \node[white, font=\tiny] at (0, 0.4) {TO-92};
            \node[below] at (0, -1.6) {Pequeña Señal};
        \end{scope}

        % TO-5 Package (Metal Can)
        \begin{scope}[shift={(3,0)}]
            % Leads
            \draw[gray, thick] (-0.2, -1.5) -- (-0.2, 0);
            \draw[gray, thick] (0.2, -1.5) -- (0.2, 0);
            \draw[gray, thick] (0, -1.5) -- (0, -0.2); % Middle lead shorter/hidden typically, but drawn for clarity

            % Body (Cylinder)
            \draw[fill=gray!40] (0,0) circle (0.5); % Top view effectively or rim
            \shade[left color=gray!30, right color=gray!60] (-0.5,0) rectangle (0.5, 0.8);
            \draw[gray] (-0.5,0) -- (-0.5,0.8);
            \draw[gray] (0.5,0) -- (0.5,0.8);
            \draw[fill=gray!20] (0,0.8) ellipse (0.5 and 0.1);

            % Tab
            \draw[fill=gray!50] (-0.6, 0.1) rectangle (-0.5, 0.3);

            \node[font=\tiny] at (0, 0.4) {TO-5};
            \node[below] at (0, -1.6) {Metálico};
        \end{scope}

        % TO-220 Package (Power)
        \begin{scope}[shift={(6,0.3)}]
            % Leads
            \draw[gray, thick] (-0.3, -1.8) -- (-0.3, -0.5);
            \draw[gray, thick] (0, -1.8) -- (0, -0.5);
            \draw[gray, thick] (0.3, -1.8) -- (0.3, -0.5);

            % Metal Tab
            \draw[fill=gray!40] (-0.5, 0.5) rectangle (0.5, 1.5);
            \draw[fill=white] (0, 1.2) circle (0.15); % Mounting hole

            % Plastic Body
            \draw[fill=black!80] (-0.5, -0.5) rectangle (0.5, 0.5);

            \node[white, font=\tiny] at (0, 0) {TO-220};
            \node[below] at (0, -1.9) {Potencia};
        \end{scope}
    \end{tikzpicture}
    \caption{Algunas formas de transistor (TO-92, TO-5, TO-220).}
    \label{fig:11-1}
\end{figure}

Algunos están montados en zócalo, como los antiguos tubos de vacío. Los zócalos se adaptan a la base de sustentación del
transistor. Algunos transistores tienen conductores flexibles para soldarlos directamente en el circuito.

La fig.~\ref{fig:11-2} es una vista en corte parcial de un transistor muy antiguo llamado de puntas
de contacto. Esta figura muestra un pequeño cristal de germanio y dos hilos delgados que hacen contacto
con el cristal. Estos elementos están encapsulados en una envolvente a prueba de humedad.

\begin{figure}[h]
    \centering
    \begin{tikzpicture}[scale=1.5]
        % Case Outline (Cutaway)
        \draw[thick, fill=gray!10] (-1.5,0) -- (-1.5,3) arc(180:0:1.5) -- (1.5,0) -- cycle;

        % Cutaway Window
        \draw[fill=white, rounded corners] (-1, 0.5) rectangle (1, 2.8);

        % Internal Structure
        % Base Support (Metal)
        \draw[fill=gray!40] (-0.8, 0.5) rectangle (0.8, 0.8);
        \node[font=\tiny, below] at (0, 0.5) {Base de Montaje};

        % Germanium Crystal (N-type typically)
        \draw[fill=black!70] (-0.4, 0.8) rectangle (0.4, 1.0);
        \node[right, font=\tiny] at (0.4, 0.9) {Cristal Germanio};

        % Whiskers
        \draw[thick] (-0.5, 2.5) .. controls (-0.5, 1.5) and (-0.1, 1.5) .. (-0.1, 1.0);
        \draw[thick] (0.5, 2.5) .. controls (0.5, 1.5) and (0.1, 1.5) .. (0.1, 1.0);

        % Whisker Supports / Leads
        \draw[thick, gray] (-0.5, 2.5) -- (-0.5, 3.2);
        \draw[thick, gray] (0.5, 2.5) -- (0.5, 3.2);

        % Labels
        \node[font=\tiny] at (0, 1.8) {Puntas de Contacto};
        \draw[->, thin] (0, 1.7) -- (-0.15, 1.2);
        \draw[->, thin] (0, 1.7) -- (0.15, 1.2);

    \end{tikzpicture}
    \caption{Vista en corte parcial de un transistor con puntas de contacto.}
    \label{fig:11-2}
\end{figure}

Los conductores terminales de estos elementos salen a través de la envolvente para establecer las conexiones
en el circuito. Los transistores de puntas de contacto están actualmente obsoletos.

La fig.~\ref{fig:11-3} es una vista en corte de un transistor de unión o contacto de aleación, y muestra una
delgada pastilla de cristal semiconductor. En cada cara de la base están aleados gránulos de indio. Los dos
gránulos y el cristal semiconductor constituyen los elementos de este transistor. Los elementos están encerrados
en una envolvente herméticamente sellada y los conductores terminales salen de la envolvente para conectarlos
en el circuito.

\begin{figure}[h]
    \centering
    \begin{tikzpicture}[scale=1.5]
        % Base Wafer (N-Type Germanium)
        \draw[fill=gray!20] (-2, 0) rectangle (2, 0.4);
        \node at (0, 0.2) {\tiny N-Type};
        \node[left, font=\tiny] at (-2.5, 0.2) {Base (Germanio)};

        % Emitter Pellet (Indium, P-Type) - Top
        \draw[fill=gray!60] (-0.5, 0.4) arc (180:0:0.5) -- cycle;
        \node[above, font=\tiny] at (0.6, 0.9) {Emisor (Indio)};

        % Collector Pellet (Indium, P-Type) - Bottom (Usually larger)
        \draw[fill=gray!60] (0.7, 0) arc (0:-180:0.7) -- cycle;
        \node[below, font=\tiny] at (0.7, -0.7) {Colector (Indio)};

        % Leads
        % Emitter Lead
        \draw[thick] (0, 0.9) -- (0, 1.5);
        \node[above] at (0, 1.5) {E};

        % Collector Lead
        \draw[thick] (0, -0.7) -- (0, -1.5);
        \node[below] at (0, -1.5) {C};

        % Base Lead (Connection to the wafer)
        \draw[thick] (-2, 0.2) -- (-2.5, 0.2);
        \draw[thick] (-2.5, -1.0) -- (-2.5, 1.0); % Base tab/ring representation
        \node[below] at (-3, 0.2) {B};

        % Junctions
        \draw[red, dashed] (-0.5, 0.4) -- (0.5, 0.4); % E-B Junction area
        \draw[red, dashed] (-0.7, 0) -- (0.7, 0);   % C-B Junction area

    \end{tikzpicture}
    \caption{Vista en corte de un transistor de unión o contacto de aleación .}
    \label{fig:11-3}
\end{figure}

\subsection{Transistores de unión}

Los transistores son realmente una extensión del diodo semiconductor. El transistor PNP ilustrado en
la fig.~\ref{fig:11-4} es un ejemplo.

\begin{figure}[h]
    \centering
    % TikZ replacement for Fig 11-4
    \begin{tikzpicture}[scale=0.8]
        % P-N-P Block
        % Emitter (P)
        \draw[fill=cyan!20] (0,0) rectangle (2,4);
        \node at (1,2) {\Large P};
        \node[left] at (-1,2) {\textbf{EMISOR}};

        % Base (N) - Thin
        \draw[fill=red!20] (2,0) rectangle (2.5,4);
        \node at (2.25,2) {\small N};
        \node[above] at (2.25,4) {\textbf{BASE}};

        % Collector (P)
        \draw[fill=cyan!20] (2.5,0) rectangle (4.5,4);
        \node at (3.5,2) {\Large P};
        \node[right] at (5.5,2) {\textbf{COLECTOR}};

        % External Box
        \draw[thick] (0,0) rectangle (4.5,4);

        % Leads
        \draw[thick] (-1,2) -- (0,2); % Emitter Lead
        \draw[thick] (2.25,0) -- (2.25,-1); % Base Lead
        \draw[thick] (4.5,2) -- (5.5,2); % Collector Lead
    \end{tikzpicture}
    \caption{Transistor PNP de unión.}
    \label{fig:11-4}
\end{figure}

Este transistor de unión está formado por un emparedado de una tira muy delgada de silicio de tipo N
entre dos tiras «anchas» de silicio de tipo P. De las placas metálicas individuales salen tres conductores
que hacen contacto con los cristales semiconductores respectivos. Todo el conjunto está alojado en una
cápsula o envolvente hermética a la humedad. La geometría de los transistores de unión difiere de la
representada en la fig.~\ref{fig:11-4}, dependiendo del método de construcción de la unión. Por otra
parte, las características de los diversos tipos de transistores de unión dependen del método de fabricación.

En la fig.~\ref{fig:11-4} la lámina o pastilla P de la izquierda se designa «emisor», la lámina N central
se designa «base» y la lámina P de la derecha se designa «colector». La base tiene aproximadamente un espesor
de 1 mil (\SI{0.001}{in}, o sea \SI{0.025}{\mm}). En lo que respecta a la polarización, este transistor
puede ser considerado como conjunto de dos diodos. El emisor y la base constituyen un diodo, el colector y
la base constituyen el otro diodo. Cuando se le utiliza como amplificador, el transistor está polarizado
como en la fig.~\ref{fig:11-5}. El diodo emisor-base está polarizado en sentido directo o de baja
resistencia por \( V_{EE} \). El diodo colector-base está polarizado en sentido inverso o de alta resistencia
por \( V_{CC} \).

\begin{figure}[h]
    \centering
    % TikZ replacement for Fig 11-5
    \begin{circuitikz}[scale=0.8]
        % P-N-P Block
        \draw[fill=cyan!20] (0,0) rectangle (2,2); \node at (1,1) {P}; % Emitter
        \draw[fill=red!20] (2,0) rectangle (2.5,2); \node at (2.25,1) {\scriptsize N}; % Base
        \draw[fill=cyan!20] (2.5,0) rectangle (4.5,2); \node at (3.5,1) {P}; % Collector

        % Leads and Circuit
        \draw (-1,1) -- (0,1); % Emitter Lead
        \draw (4.5,1) -- (5.5,1); % Collector Lead
        \draw (2.25,0) -- (2.25,-0.5); % Base Lead (internal start)

        % Circuit Loop Left (Emitter-Base)
        \draw (-1,1) -- (-1,-2)
        to[battery1, l=$V_{EE}$] (2.25,-2)
        -- (2.25, -0.5);

        % Circuit Loop Right (Collector-Base)
        \draw (5.5,1) -- (5.5,-2)
        to[battery1, l=$V_{CC}$, invert] (2.25,-2);

        % Labels
        \node[above] at (-1,1) {Emisor};
        \node[above] at (2.25,2) {Base};
        \node[above] at (5.6,1) {Colector};
    \end{circuitikz}
    \caption{Polarización de un transistor PNP.}
    \label{fig:11-5}
\end{figure}

Los huecos son los portadores mayoritarios de corriente en el diodo emisor-base y emanan del emisor
del tipo P. Si no estuviese presente el colector, el flujo de corriente en la sección emisor-base
tendría lugar de la manera descrita en la Práctica 1. Sin embargo, la presencia del colector de tipo
P y su conexión al terminal negativo de la batería de colector \( V_{CC} \) altera radicalmente la
trayectoria o camino del flujo de la corriente de huecos. Sólo un pequeño porcentaje de los huecos
emitidos por el emisor se combinan con los electrones libres en la base. Los otros huecos
(aproximadamente 95\%) atraviesan el cristal muy delgado de la base y son atraídos por el terminal
negativo de la batería en el colector. El circuito emisor-colector está completado externamente a
través de las dos baterías \( V_{EE} \) y \( V_{CC} \) conectadas en serie aditiva. De esta descripción
del flujo de corriente en un transistor se deduce (fig. 11-5) que el emisor es la fuente de los portadores
de corriente, que la corriente emisor-base es muy pequeña y que la corriente emisor-colector es muy intensa.
También se deduce que los cambios de la polarización emisor-base darán por resultado variaciones de la
corriente de emisor. Así pues, un aumento de la polarización directa dará por resultado un aumento de
la corriente de emisor y por tanto de la corriente de colector. La corriente de base aumentará o
disminuirá muy poco cuando aumente o disminuya la corriente de emisor. Es pues evidente que la corriente
de colector se puede controlar fácilmente variando la polarización emisor-base. Las designaciones
«emisor» y «colector» se pueden ahora asociar fácilmente con sus funciones.

También puede estar fabricado un transistor de unión con la configuración NPN como en la
fig.~\ref{fig:11-6}. Lo mismo que en el caso precedente, la polarización emisor-base debe tener sentido
directo y la de colector-base sentido inverso. A causa de que aquí se utiliza un cristal de tipo N como
emisor y uno de tipo P como base, las polaridades de la batería deben estar invertidas, con respecto
a la polarización de un transistor PNP. En el transistor NPN los electrones son los portadores mayoritarios
de corriente. En el transistor PNP los portadores mayoritarios son los huecos.

\begin{figure}[h]
    \centering
    % TikZ replacement for Fig 11-6
    \begin{circuitikz}[scale=0.8]
        % N-P-N Block
        \draw[fill=red!20] (0,0) rectangle (2,2); \node at (1,1) {N}; % Emitter
        \draw[fill=cyan!20] (2,0) rectangle (2.5,2); \node at (2.25,1) {\scriptsize P}; % Base
        \draw[fill=red!20] (2.5,0) rectangle (4.5,2); \node at (3.5,1) {N}; % Collector

        % Leads and Circuit
        \draw (-1,1) -- (0,1); % Emitter Lead
        \draw (4.5,1) -- (5.5,1); % Collector Lead
        \draw (2.25,0) -- (2.25,-0.5); % Base Lead

        % Circuit Loop Left (Emitter-Base)
        \draw (-1,1) -- (-1,-2) to[battery1, l=$V_{EE}$, invert] (2.25,-2) -- (2.25,-0.5);

        % Circuit Loop Right (Collector-Base)
        \draw (5.5,1) -- (5.5,-2) to[battery1, l=$V_{CC}$] (2.25,-2);

        % Labels
        \node[above] at (-1,1) {Emisor};
        \node[above] at (2.25,2) {Base};
        \node[above] at (5.6,1) {Colector};
    \end{circuitikz}
    \caption{Polarización de un transistor NPN.}
    \label{fig:11-6}
\end{figure}

Un esquema sencillo para recordar las conexiones de los polos o terminales de la batería aplicados
a un transistor es el siguiente: las polaridades de la batería aplicadas al emisor y la base tienen
por iniciales respectivamente las dos primeras letras de la designación del tipo. Así pues, en el tipo
PNP el emisor recibe la polaridad P (positiva) y la base recibe la polaridad N (negativa). Análogamente,
en el tipo NPN, el emisor recibe la polaridad N (negativa) y la base la polaridad P (positiva).
El colector recibe una polaridad de batería opuesta a la designación de su tipo de impureza. El colector
de tipo P recibe pues la polaridad N (negativa) de la batería y el colector de tipo N recibe la
polaridad P (positiva).

La fig.~\ref{fig:11-7} es un diagrama que indica el sentido del flujo de electrones en el circuito
exterior de un transistor PNP. Se puede considerar que en el circuito exterior se establece una corriente
o flujo de electrones y que dentro del cristal del tipo P hay un flujo de huecos. La corriente designada
por \( I_{CBO} \) es muy pequeña y aquí no trataremos de ella. La corriente en el circuito de emisor es
la «total» y se divide en la corriente de base y la corriente de colector. Es importante señalar que la
corriente emisor-base (determinada por la tensión de polarización base-emisor) es realmente la causa
de la corriente emisor-colector. Como veremos más adelante, la corriente emisor-base también controla
la corriente emisor-colector.

\begin{figure}[h]
    \centering
    % TikZ replacement for Fig 11-7
    \begin{circuitikz}[scale=0.8]
        % P-N-P Block
        \draw[fill=cyan!20] (0,0) rectangle (2,2); \node at (1,1) {P}; % Emitter
        \draw[fill=red!20] (2,0) rectangle (2.5,2); \node at (2.25,1) {\scriptsize N}; % Base
        \draw[fill=cyan!20] (2.5,0) rectangle (4.5,2); \node at (3.5,1) {P}; % Collector

        % Leads
        \draw (-1,1) -- (0,1); % Emitter Lead
        \draw (4.5,1) -- (5.5,1); % Collector Lead
        \draw (2.25,0) -- (2.25,-0.5); % Base Lead (internal start)

        % Circuit Loop Left (Emitter-Base)
        \draw (-1,1) -- (-1,-2)
        to[battery1, l=$V_{EE}$] (2.25,-2)
        -- (2.25, -0.5);

        % Circuit Loop Right (Collector-Base)
        \draw (5.5,1) -- (5.5,-2)
        to[battery1, l=$V_{CC}$, invert] (2.25,-2);

        % Electron Flow Arrows
        \draw[->, thick, blue] (2.25, -1.5) -- (2.25, -0.8) node[midway, right] {$e^-$}; % Into Base (from VEE)
        \draw[->, thick, blue] (-1, 0.5) -- (-1, -0.5) node[midway, left] {$e^-$}; % Out of Emitter (to VEE)

        \draw[->, thick, blue] (5.5, -0.5) -- (5.5, 0.5) node[midway, right] {$e^-$}; % Into Collector (from VCC)
        \draw[->, thick, blue] (2.25, -0.8) -- (2.25, -1.5); % Out of Base (to VCC)

        % Labels
        \node[above] at (-1,1) {Emisor};
        \node[above] at (2.25,2) {Base};
        \node[above] at (5.6,1) {Colector};
    \end{circuitikz}
    \caption{Flujo de electrones en el circuito externo de un transistor PNP.}
    \label{fig:11-7}
\end{figure}

La fig.~\ref{fig:11-8} es un diagrama simplificado que indica el sentido del flujo de electrones en el
circuito exterior de un transistor NPN. Comparado con el de la fig.~\ref{fig:11-7} se observa que el sentido del
flujo de electrones en el circuito exterior de un transistor NPN es opuesto al del tipo PNP.

\begin{figure}[h]
    \centering
    % TikZ replacement for Fig 11-8
    \begin{circuitikz}[scale=0.8]
        % N-P-N Block
        \draw[fill=red!20] (0,0) rectangle (2,2); \node at (1,1) {N}; % Emitter
        \draw[fill=cyan!20] (2,0) rectangle (2.5,2); \node at (2.25,1) {\scriptsize P}; % Base
        \draw[fill=red!20] (2.5,0) rectangle (4.5,2); \node at (3.5,1) {N}; % Collector

        % Leads
        \draw (-1,1) -- (0,1); % Emitter Lead
        \draw (4.5,1) -- (5.5,1); % Collector Lead
        \draw (2.25,0) -- (2.25,-0.5); % Base Lead

        % Circuit Loop Left (Emitter-Base)
        \draw (-1,1) -- (-1,-2)
        to[battery1, l=$V_{EE}$, invert] (2.25,-2)
        -- (2.25, -0.5);

        % Circuit Loop Right (Collector-Base)
        \draw (5.5,1) -- (5.5,-2)
        to[battery1, l=$V_{CC}$] (2.25,-2);

        % Electron Flow Arrows
        \draw[->, thick, blue] (-1, -0.5) -- (-1, 0.5) node[midway, left] {$e^-$};

        \draw[->, thick, blue] (5.5, 0.5) -- (5.5, -0.5) node[midway, right] {$e^-$};

        \draw[->, thick, blue] (2.25, -0.8) -- (2.25, -1.5); % To V_EE (+)
        \draw[->, thick, blue] (2.25, -1.5) -- (2.25, -0.8); % From V_CC (-)

        % Labels
        \node[above] at (-1,1) {Emisor};
        \node[above] at (2.25,2) {Base};
        \node[above] at (5.6,1) {Colector};
    \end{circuitikz}
    \caption{Flujo de electrones en el circuito externo de un transistor NPN.}
    \label{fig:11-8}
\end{figure}

\subsection{\texorpdfstring{\( I_{CBO} \)}{I\_CBO} y el embalamiento térmico}

Una característica importante del circuito colector-base es \( I_{CO} \) (llamada también
\( I_{CBO} \)), corriente de colector que se establece estando la unión colector-base polarizada
en sentido inverso y la de emisor-base en cortocircuito. Esta corriente de fuga (\( I_{CBO} \)) es
debida a los portadores minoritarios en la base y en el colector. \( I_{CBO} \) es del orden de
algunos microamperios para el germanio y de algunos nanoamperios para el silicio, y aumenta con
la temperatura.

Un factor importante que afecta al funcionamiento de un transistor es su temperatura de trabajo.
El transistor es muy sensible a los cambios de temperatura. Cuando ésta aumenta, la corriente también
aumenta y ésta a su vez hace que el calor aumente. Si esta reacción en cadena, que se llama
«embalamiento», no se interrumpe, puede conducir a la destrucción completa del transistor a causa
del exceso de calor. En el diseño de los circuitos de transistor se evita el embalamiento térmico
mediante las disposiciones de realimentación negativa que estudiaremos. El margen normal de temperatura
dentro del cual puede funcionar un transistor en condiciones de seguridad está especificado por el
fabricante. Los transistores de silicio toleran más el calor que los de germanio, y su margen
térmico de funcionamiento es mucho más amplio.

\begin{figure}[h]
    \centering
    % TikZ replacement for Fig 11-9
    \begin{circuitikz}[scale=1.2]
        \draw (0,0) node[pnp, text width=1cm] (Q1) {} node[below=0.5cm of Q1] {(a) PNP};
        \node[above] at (Q1.E) {E};
        \node[below] at (Q1.C) {C};
        \node[left] at (Q1.B) {B};

        \draw (4,0) node[npn, text width=1cm] (Q2) {} node[below=0.5cm of Q2] {(b) NPN};
        \node[below] at (Q2.E) {E};
        \node[above] at (Q2.C) {C};
        \node[left] at (Q2.B) {B};
    \end{circuitikz}
    \caption{Símbolo esquemático de transistores (a) PNP y (b) NPN.}
    \label{fig:11-9}
\end{figure}

\subsection{Símbolo de transistor, zócalo y montura}

Los símbolos de transistores PNP y NPN están indicados en las fig.~\ref{fig:11-9} a y b, donde el
elemento correspondiente a la flecha es el emisor y el otro simétrico es el colector. El transistor
PNP se caracteriza porque la flecha del emisor está dirigida hacia la base mientras que en el tipo
NPN la flecha está dirigida en sentido contrario a la base.

Hay dos disposiciones generales de terminales del transistor. En la primera se emplean conductores
largos y flexibles.
Ejemplo es el transistor representado en la fig.~\ref{fig:11-10}a. Los tres conductores están alineados
y corresponden a emisor, base y colector, en este orden. La separación del terminal de colector es mayor
que la de los otros dos elementos. En una variante de este sistema se utilizan también tres conductores
alineados y el conductor se identifica por un punto rojo pintado cerca de la caja, como en la
fig.~\ref{fig:11-10}b. El método usual para conectar los transistores de las fig.~\ref{fig:11-10} a y b en el
circuito es soldar los conductores flexibles a los terminales correspondientes del circuito.

\begin{figure}[h]
    \centering
    \begin{tikzpicture}[scale=1.2]
        % (a) Identification by lead spacing
        \begin{scope}[shift={(0,0)}]
            % Case Outline (Bottom View - Ellipse/Circle)
            \draw[thick, fill=gray!20] (0,0) ellipse (1.2 and 1);

            % Leads
            \node[circle, fill=black, inner sep=1.5pt, label=below:E] (E) at (-0.6, -0.2) {};
            \node[circle, fill=black, inner sep=1.5pt, label=below:B] (B) at (0, -0.2) {};
            \node[circle, fill=black, inner sep=1.5pt, label=below:C] (C) at (0.8, -0.2) {};

            % Spacing indications
            \draw[<->, blue, thin] (E) -- (B) node[midway, above=-2pt] {\tiny $x$};
            \draw[<->, blue, thin] (B) -- (C) node[midway, above=-2pt] {\tiny $>x$};

            \node[below=1cm] at (0,-0.5) {(a)};
        \end{scope}

        % (b) Identification by colored dot
        \begin{scope}[shift={(4,0)}]
            % Case Outline
            \draw[thick, fill=gray!20] (0,0) ellipse (1.2 and 1);

            % Leads (Equal spacing)
            \node[circle, fill=black, inner sep=1.5pt, label=below:E] (E2) at (-0.6, -0.2) {};
            \node[circle, fill=black, inner sep=1.5pt, label=below:B] (B2) at (0, -0.2) {};
            \node[circle, fill=black, inner sep=1.5pt, label=below:C] (C2) at (0.6, -0.2) {};

            % Red Dot
            \fill[red] (0.9, 0) circle (0.15);
            \node[red, right] at (1.2, 0) {\tiny Punto Rojo};

            \node[below=1cm] at (0,-0.5) {(b)};

            % Lead indication for Red Dot (often C)
            \draw[->, red, dashed] (0.9, -0.1) -- (C2);
        \end{scope}

    \end{tikzpicture}
    \caption{Métodos para identificar los terminales de los transistores.}
    \label{fig:11-10}
\end{figure}

En el segundo método se utilizan patillas rígidas a las que están conectados los conductores interiores,
como en la fig.~\ref{fig:11-11}a. Este es un tipo enchufable de transistor y requiere un zócalo o enchufe
para conectarlo. Otra disposición de terminales permite elegir el primero o el segundo método de montaje.
En este caso el transistor está provisto de conductores flexibles soldados a patillas rígidas (fig. 11-2).
Los conductores flexibles se emplean cuando haya que soldar el transistor en el circuito. Si se quiere
emplear el zócalo, se quitan fácilmente los conductores flexibles con alicates. También se usa un tipo
triangular de zócalo (figs.~\ref{fig:11-11}b y c). Otras disposiciones típicas de terminales están
representadas en la fig.~\ref{fig:11-12}.

\begin{figure}[h]
    \centering
    \begin{tikzpicture}[scale=1.2]
        % (a) Plug-in Transistor (Side View)
        \begin{scope}[shift={(0,0)}]
            % Pins
            \draw[thick, gray] (-0.3, -1) -- (-0.3, 0);
            \draw[thick, gray] (0, -1) -- (0, 0);
            \draw[thick, gray] (0.3, -1) -- (0.3, 0);

            % Body
            \draw[fill=gray!30] (-0.5, 0) rectangle (0.5, 0.8);
            \draw[fill=gray!10] (0, 0.8) ellipse (0.5 and 0.1);
            \node at (0, 0.4) {\tiny Transistor};
            \node[below=1.2cm] at (0,0) {(a) Transistor};
        \end{scope}

        % (b) Base Diagram (Bottom View)
        \begin{scope}[shift={(3,0)}]
            \draw[thick] (0,0.1) circle (1.1);
            \node at (0,1.4) {\tiny Vista Inferior};

            % Triangular Pin Layout
            % Isosceles triangle.
            \node[circle, fill=black, inner sep=2pt, label=left:E] (E) at (-0.4, -0.2) {};
            \node[circle, fill=black, inner sep=2pt, label=right:C] (C) at (0.4, -0.2) {};
            \node[circle, fill=black, inner sep=2pt, label=above:B] (B) at (0, 0.5) {};

            \draw[dashed] (E) -- (B) -- (C) -- cycle;
            \node[below=1.2cm] at (0,0) {(b) Base};
        \end{scope}

        % (c) Socket (Top View)
        \begin{scope}[shift={(6,0)}]
            \draw[fill=black!10] (-0.8, -0.8) rectangle (0.8, 0.8);
            \node at (0,1) {\tiny Zócalo};

            % Holes
            \fill[white, draw=black] (-0.4, -0.2) circle (0.1);
            \fill[white, draw=black] (0.4, -0.2) circle (0.1);
            \fill[white, draw=black] (0, 0.5) circle (0.1);

            \node[below=1.2cm] at (0,0) {(c) Zócalo};
        \end{scope}
    \end{tikzpicture}
    \caption{(a) Transistor enchufable. Identificación de la base triangular;
        (b) diagrama de la base (vista inferior);
        (c) zócalo (CBS).}
    \label{fig:11-11}
\end{figure}

\begin{figure}[h]
    \centering
    \begin{tikzpicture}[scale=1.2]
        % (a) TO-3 (Power) - Bottom View
        \begin{scope}[shift={(0,0)}]
            % Diamond Shape
            \draw[thick, fill=gray!30] (0, 1) -- (1.5, 0) -- (0, -1) -- (-1.5, 0) -- cycle;
            \draw[fill=white] (-1.2, 0) circle (0.1); % Mounting hole
            \draw[fill=white] (1.2, 0) circle (0.1); % Mounting hole

            % Pins (Offset from center usually)
            \node[circle, fill=black, inner sep=2pt, label=below:E] (E) at (-0.4, 0.3) {};
            \node[circle, fill=black, inner sep=2pt, label=below:B] (B) at (0.4, 0.3) {};
            \node at (0, -0.5) {C (Caja)};

            \node[below=1.8cm] at (0,0) {(a) TO-3 (Vista Inf.)};
        \end{scope}

        % (b) TO-220 - Front View
        \begin{scope}[shift={(3.5,0)}]
            % Heatsink Tab
            \draw[fill=gray!40] (-0.6, 0.5) rectangle (0.6, 1.2);
            \draw[fill=white] (0, 0.9) circle (0.1); % Hole

            % Plastic Body
            \draw[fill=black!80] (-0.6, -0.5) rectangle (0.6, 0.5);

            % Pins
            \draw[thick, gray] (-0.4, -0.5) -- (-0.4, -1.2) node[below, black] {\tiny B};
            \draw[thick, gray] (0, -0.5) -- (0, -1.2) node[below, black] {\tiny C};
            \draw[thick, gray] (0.4, -0.5) -- (0.4, -1.2) node[below, black] {\tiny E};

            \node[below=1.8cm] at (0,0) {(b) TO-220 (Frente)};
        \end{scope}

        % (c) TO-92 - Bottom View
        \begin{scope}[shift={(7,0)}]
            % D-Shape
            \draw[thick, fill=gray!20] (0.6, 0.0) arc (0:180:0.6) -- cycle;

            % Pins
            \node[circle, fill=black, inner sep=1.5pt, label={[yshift=-2mm]below:E}] (E92) at (-0.4, 0.2) {};
            \node[circle, fill=black, inner sep=1.5pt, label={[yshift=-2mm]below:B}] (B92) at (0, 0.1) {};
            \node[circle, fill=black, inner sep=1.5pt, label={[yshift=-2mm]below:C}] (C92) at (0.4, 0.2) {};

            \node[below=1.8cm] at (0,0) {(c) TO-92 (Vista Inf.)};
        \end{scope}

    \end{tikzpicture}
    \caption{Otras disposiciones típicas de terminales (TO-3, TO-220, TO-92).}
    \label{fig:11-12}
\end{figure}

Los transistores se especifican según su aptitud para disipar potencias. Un transistor empleado
como amplificador de audio de bajo nivel puede tener una potencia nominal baja, del orden de
\SI{50}{\milli\watt}. Un transistor utilizado como amplificador de salida debe tener una
especificación en vatios más alta. Las cajas o cápsulas de los transistores de potencia están
diseñadas especialmente para facilitar su refrigeración. Por ejemplo, algunos transistores de
potencia usan aletas radiales para conducir el calor y disiparlo (fig. 11-1). Otros tipos usan
un casco o caja metálica que se monta en el chasis metálico del equipo en que haya de emplearse
(fig. 11-3). Este transistor tiene el colector conectado a la caja del transistor. Entonces el
chasis conduce el calor y lo disipa. Este tipo de transistores es de potencia nominal más alta
cuando está fijado al chasis metálico y de potencia nominal más baja cuando se monta fuera del
chasis. Una buena técnica de diseño requiere que los transistores estén montados en la parte más
fría del chasis. Algunas veces hay aletas instaladas dentro del equipo para refrigerar los
transistores y otros componentes. En la lista siguiente se numeran los símbolos utilizados para
denotar los parámetros del transistor.

\begin{table}[H]
    \centering
    \begin{tabular}{cp{12cm}}
        \toprule
        Símbolo      & Definición                                                                    \\
        \midrule
        \( I_C \)    & Corriente de colector                                                         \\
        \( I_E \)    & Corriente de emisor                                                           \\
        \( I_B \)    & Corriente de base                                                             \\
        \( V_C \)    & Tensión en el colector                                                        \\
        \( V_E \)    & Tensión en el emisor                                                          \\
        \( V_B \)    & Tensión en la base                                                            \\
        \( V_{CC} \) & Tensión de alimentación del colector                                          \\
        \( V_{EE} \) & Tensión de alimentación del emisor                                            \\
        \(V_{BB} \)  & Tensión de alimentación de la base                                            \\
        \( I_{CB} \) & Corriente de colector-base (el segundo subindice se puede utilizar
        en cualquiera de los símbolos anteriores para evitar confusiones)                            \\
        \(V_{KJ} \)  & Tensión del circuito entre elementos, por ejemplo, entre los elementos K y J. \\
        \(V_{CB} \)  & Tensión colector-base                                                         \\
        \bottomrule
    \end{tabular}
    \caption{Símbolos utilizados para denotar los parámetros del transistor.}
    \label{tab:symbols}
\end{table}

\section{Reglas para Trabajar con Circuitos de Transistor y Efectuar
  Mediciones en Ellos.}
\subsection{Herramientas especiales}
Se requiere el uso de herramientas especiales para trabajar en circuitos
de transistor a causa del pequeño tamaño de estos y de sus componentes
asociados. Por otra parte, como los equipos transistorizados suelen estar
miniaturizados, las herramientas ordinarias serían inaplicables. He aquí una
lista de herramientas especiales:
\begin{itemize}
    \item Pinzas miniaturas.
    \item Alicates de punta miniatura.
    \item Surtido de tenacillas del tipo de servicio.
    \item Un soldador de \SI{25}{\watt} del tipo de lápiz.
    \item Un soldador auxiliar.
    \item Una lente de relojero o una lupa con pié.
\end{itemize}

\subsection{Precauciones para la manipulación de los transistores.}
Aunque son dispositivos robustos, los transistores se deterioran
fácilmente si se les maneja inapropiadamente. Por esto hay que poner gran
cuidado cuando se trabaja con transistores de baja potencia que tienen
conductores flexibles, ya que estos son frágiles y se rompen con facilidad.
Las mismas precauciones son aplicables a los componentes de los equipos
transistorizados. Estos componentes han sido miniaturizados y sus conductores
se pueden romper fácilmente. Los transistores y los componentes asociados
empleados en esta práctica y las siguientes se deben montar en paneles de
ensayo. Los transistores se pueden destruir casi instantáneamente a diferencia
de los tubos de vacío que pueden tolerar moderadas sobrecargas durante
largos períodos de tiempo. Un cortocircuito entre base y colector en un
circuito en funcionamiento destruirá casi siempre al transistor. Asimismo,
un cortocircuito en una sonda de medición o una herramienta entre un par
de terminales puede destruir un transistor.

Hay que poner un cuidado extremado cuando se insertan transistores en un
circuito, o cuando se desmontan, estando aplicada la potencia. Lo mejor
será siempre desconectar la potencia antes de hacer estas operaciones.

\subsection{Otras precauciones.}
Polarización del colector y valores de tensión. Asegurarse de que la tensión
de polarización del colector es correcta antes de aplicar la potencia. El
colector debe tener polarización inversa. Por otra parte, la tensión en
él y en el emisor no debe exceder los valores especificados.
Por consiguiente habrá que medir estos valores y ajustarlos al valor correcto antes de
aplicar la potencia al circuito.

\subsection{Comprobación de las conexiones del circuito.}
Se deben comprobar todas las conexiones a la vista del esquema antes de aplicar la potencia.
Los transistores no deben ser conectados a una fuente de tensión sin una resistencia limitadora
en el circuito.

\subsection{Operación de soldadura de los terminales de los transistores.}
Debe ser todo lo rápida posible; a este fin son recomendables los
soldadores de poco vatiaje (\SI{25}{\watt}). Los terminales se deben dejar todo lo
largos que sea posible, pero teniendo en cuenta las consideraciones de
diseño del circuito para reducir la transferencia de calor. Esto mismo
es aplicable a los diodos de germanio. Durante la operación de soldadura
se improvisa un absorbedor y radiador de calor manteniendo sujetos con
los alicates los terminales del transistor entre el cuerpo de éste y el
punto de aplicación del calor, como en la fig.~\ref{fig:11-13}.

\begin{figure}[h]
    \centering
    \begin{tikzpicture}[scale=1.2]
        % Transistor Body
        \draw[fill=gray!30] (-0.5, 2) rectangle (0.5, 2.8);
        \node at (0, 2.4) {\tiny Transistor};

        % Leads
        \draw[thick, gray] (-0.3, 2) -- (-0.3, 0);
        \draw[thick, gray] (0, 2) -- (0, 0);
        \draw[thick, gray] (0.3, 2) -- (0.3, 0);

        % Pliers (Heat Sink) gripping the right lead (0.3, y)
        % Pliers stylized as two jaws clamping
        \draw[fill=gray!60, draw=black] (0.35, 1.0) -- (1.5, 1.2) -- (1.5, 0.8) -- (0.35, 1.0); % Right Jaw handle part
        \draw[fill=gray!60, draw=black] (0.25, 1.0) -- (-0.9, 1.2) -- (-0.9, 0.8) -- (0.25, 1.0); % Left Jaw (behind/front illusion)

        % Better Pliers Representation: Side view clamping the wire
        % Jaws
        \draw[fill=gray!80] (0.1, 1.1) rectangle (0.5, 1.3); % Top Jaw
        \draw[fill=gray!80] (0.1, 0.7) rectangle (0.5, 0.9); % Bottom Jaw
        % Handles extending out
        \draw[gray, thick] (0.5, 1.2) -- (2, 1.5);
        \draw[gray, thick] (0.5, 0.8) -- (2, 0.5);

        \node[right, align=left, font=\small] at (2, 1) {Alicates\\(Radiador de calor)};
        \draw[->, thin] (1.9, 1) -- (0.6, 1);

        % Soldering Iron Tip at the bottom of the lead
        \draw[fill=orange!40] (0.3, -0.1) -- (0.5, -0.6) -- (0.1, -0.6) -- cycle; % Tip
        \draw[fill=gray!50] (0.1, -0.6) rectangle (0.5, -1.2); % Iron Body
        \node[below] at (0.3, -1.2) {Soldador};

        % Heat indication
        \draw[->, red, thick, decorate, decoration={snake, amplitude=.4mm, segment length=2mm, post length=1mm}] (0.3, -0.2) -- (0.3, 0.6); % Heat moving up
        \node[right, red] at (0.4, 0) {Calor};

        % Heat blocked/dissipated by pliers
        \draw[->, blue, thick] (0.3, 0.9) -- (0.8, 2); % Heat dissipating into pliers
        %\draw[->, blue, thick] (0.3, 0.9) -- (0.8, 1);
        \node[right, blue] at (0.8, 2) {Disipación};

    \end{tikzpicture}
    \caption{Los alicates actúan como radiador de calor cuando se sueldan los transistores en el circuito.}
    \label{fig:11-13}
\end{figure}

\subsection{Medición de tensión.}
Todo el equipo de ensayo o prueba debe estar aislado de la red o línea
de potencia. Si el equipo no está aislado, habrá que intercalar un
transformador. Cuando se efectúan mediciones de tensión en circuitos de
transistor, hay que adoptar las debidas precauciones para evitar cortocircuitos
accidentales entre terminales poco separados. Estos cortocircuitos pueden ser causa
de que se aplique una tensión incorrecta o excesiva a los elementos del transistor
y de la consiguiente destrucción de éste.

\subsection{Mediciones de resistencia.}
Los transistores se pueden estropear mediante las mediciones de
resistencia. Si se emplea un óhmetro de tipo shunt en los alcances o
escalas de baja resistencia, la corriente a que se somete al transistor
puede ser excesiva y destruirlo. Por consiguiente habrá que adoptar
precauciones cuando se efectúan las pruebas con un óhmetro en los
circuitos de transistor. Una buena regla es comprobar el potencial y la
polaridad de los terminales del óhmetro antes de probar las uniones del
transistor.

También es necesario sacar el transistor primero, si es del tipo
enchufable, cuando se hacen pruebas de resistencia de los componentes
en un circuito transistorizado. Cuando se efectúan las mediciones de
resistencia estando el transistor intercalado en el circuito, habrá que
tener en cuenta la corriente admisible de conducción en el transistor. Un
método recomendable es hacer dos series de mediciones de resistencia,
invirtiendo los conductores del óhmetro para hacer la lectura cuando sea
necesario. La lectura más alta es la más correcta a causa de que
corresponde a la polarización inversa.

Ordinariamente, los voltímetros electrónicos (EVM) contienen una función
*ohmios-baja potencia* (LPD) destinada especialmente a mediciones de
resistencia en circuitos transistorizados. La tensión que se desarrolla en
la función LPD es insuficiente para polarizar en sentido directo las
uniones del transistor. Por esto se utiliza dicha función LPD para medir
resistencias en circuitos que contienen dispositivos de estado sólido.

\subsection{Uso de los generadores de señal en los circuitos de transistor.}
Una salida excesiva del generador de señal puede destruir el transistor.
Por tanto, la salida debe ser ajustada al mínimo cuando se invierta
inicialmente la señal. Como medida adicional de seguridad, el generador no
debe ser acoplado directamente al circuito. Siempre que sea posible
conviene que el acoplamiento sea flojo.

\section{Resumen}
\begin{enumerate}
    \item Los transistores son dispositivos de estado sólido que extienden el margen y la
          aplicación de los diodos de dos elementos.
    \item El silicio es el elemento con el que están construidos la mayoría de transistores
          y otros semiconductores actualmente.
    \item Los transistores de unión o contacto están constituidos por un elemento muy delgado
          llamado base, emparedado entre dos elementos llamados emisor y colector (fig. 11-4).
    \item Los transistores son del tipo PNP o del tipo NPN. La primera letra designa el material
          de emisor, la letra central designa la base y la última letra designa el material de colector.
          Por tanto, los transistores contienen dos uniones, la unión emisor-base y la unión colector-base.
    \item Las características de los transistores de unión dependen del grado de dopado del emisor,
          de la base y del colector; dependen también de la geometría del transistor y del método de fabricación.
    \item En lo que respecta a la polarización, el transistor se puede considerar constituido por
          dos diodos, el diodo emisor-base y el diodo colector-base.
    \item En la mayoría de aplicaciones el diodo emisor-base está polarizado en sentido directo y el
          diodo colector-base lo está en sentido inverso.
    \item El emisor *inyecta* (emite) o es la fuente de los portadores de corriente en un transistor. El
          colector recibe (recoge) la mayoría de los portadores de corriente. La base controla la corriente de colector.
    \item La corriente de emisor es la corriente total, y se divide en la corriente de colector y en
          la corriente de base.
    \item La corriente que fluye dentro de un transistor está constituida por los portadores mayoritarios,
          que son electrones en el material de tipo N y huecos en el material de tipo P.
    \item Hay también portadores *minoritarios*. La pequeña corriente de portadores minoritarios en la unión
          colector-base estando el emisor abierto se llama corriente de fuga o \( I_{CBO} \).
    \item Aunque \( I_{CBO} \) es muy pequeña en los transistores de silicio, aumenta con el calor y puede
          destruir un transistor si no se controla.
    \item Los transistores son sensibles al calor. Las temperaturas a las que pueden trabajar están
          especificadas por el fabricante.
    \item Los transistores pueden ser del tipo de terminales soldados o del tipo enchufable. En el
          primer tipo se utilizan zócalos cuya forma y tamaño varía para adaptarse al transistor.
    \item Los transistores se especifican de acuerdo con su aptitud para disipar potencia, variando
          desde potencia muy baja (50 mW) hasta potencia alta (2 W o más).
    \item En los transistores de alta potencia se emplean radiadores de calor para conducir éste
          desde el transistor y disiparlo.
    \item En la operación de soldadura de los transistores se utilizan soldadores de bajo vatiaje
          para evitar que se calienten excesivamente.
    \item Los transistores se deben manejar cuidadosamente para no dañar sus frágiles conductores.
    \item Siempre se debe desconectar la alimentación de potencia antes de insertar o sacar un
          transistor en un circuito.
    \item Los transistores se pueden destruir accidentalmente cuando se miden tensiones en sus terminales.
          Hay que poner cuidado para evitar que se cortocircuiten los terminales durante las mediciones en el circuito.
    \item Los transistores se pueden deteriorar durante las pruebas de sus elementos con el óhmetro si
          éste aplica una corriente excesiva al transistor. Por tanto se deben emplear alcances bajos de
          corriente en el óhmetro.
    \item En las pruebas de resistencia de un circuito que contiene un transistor u otro dispositivo de
          estado sólido se debe utilizar la función de óhmetro de baja potencia en un EVM. Así se evita que
          los electrodos o elementos del transistor se polaricen en sentido directo y se garantiza la obtención
          de lecturas correctas de la resistencia.
\end{enumerate}

\begin{comment}
\section{Autoexamen}
\begin{enumerate}
    \item Actualmente los transistores se fabrican con ...... (germonio/silicio).
    \item Los transistores de tres elementos contienen ...... (2/3 uniones).
    \item La unión emisor-base de un transistor está polarizada ......; la unión
          colector-base está polarizada ......
    \item En un transistor PNP el emisor es ...... (positivo/negativo) con respecto a la base,
          mientras el colector es ...... (positivo/negativo) con respecto a la base.
    \item Las polaridades de batería en un transistor NPN ...... (están invertidas/son las mismas)
          en comparación con las de un transistor PNP.
    \item La cabeza de flecha del símbolo de emisor de un transistor indica el sentido de la corriente
          convencional, que es opuesto al del flujo de electrones ...... (verdad/falso).
    \item Los transistores son dispositivos con los que no es necesario poner mucho cuidado cuando se
          les maneja ...... (verdad/falso).
    \item Un punto en la caja del transistor o la separación entre los elementos es lo que se suele
          utilizar para identificar los elementos del transistor ...... (verdad/falso).
    \item Los transistores son sensibles ...... Por tanto se deben emplear ...... cuando se sueldan
          los terminales del transistor en el circuito.
    \item Se debe: ...... cuando se efectúan mediciones de tensión entre los elementos de un transistor
          incorporado en un circuito para evitar ...... del transistor.
\end{enumerate}
\end{comment}

\section{Materiales Necesarios}
\begin{itemize}
    \item Fuente de alimentación: fuentes de c.c. de 1,5 y 6 V.
    \item Equipo: 2 miliamperímetros de varios alcances (o VOM de 20.000 \(\Omega\)/V); EVM.
    \item Resistores: \( \frac{1}{2} \) W 100 y 820 \(\Omega\).
    \item Semiconductores: Transistores 2N6004 y 2N6005 montados apropiadamente o equivalente.
    \item Diversos: Potenciómetro de 2 W 2.500 \(\Omega\).
    \item Dos interruptores unipolares.
\end{itemize}

\section{Procedimiento}

\subsection{Polarización PNP}
\begin{enumerate}
    \item El alumno recibirá un transistor de unión PNP y uno NPN, correctamente identificado.
          Examinar estos transistores. Dibujar un diagrama básico de cada uno, identificando el
          transistor y cada uno de los elementos.
    \item Conectar el circuito de la fig.~\ref{fig:11-14} con el interruptor de potencia \( S_1 \) abierto.
          Ajustar \( R_2 \) para la máxima resistencia. Así se aplicará la mínima polarización al emisor
          cuando se aplique la potencia. \( R_1 \) es un resistor limitador de corriente (polarización) en el
          circuito de emisor. \( R_2 \) es un resistor limitador en el circuito de colector. Mantener la polaridad
          correcta de la batería y del medidor. Ajustar los medidores \( M_1 \) y \( M_2 \) en sus alcances altos.
          Después de aplicar la potencia, disminuir el alcance del medidor hasta que sea el adecuado para leer la corriente.
          El instructor debe aprobar las conexiones del circuito antes de aplicar la potencia.

          \begin{figure}[H]
              \centering
              \begin{circuitikz}[american, scale=1.0]
                  % PNP Transistor
                  % Rotate 90: E Left, C Right, B Bottom
                  \draw (0,0) node[pnp, rotate=90] (Q1) {};
                  \node[right] at (Q1.B) {B};
                  \node[below] at (Q1.E) {E};
                  \node[below] at (Q1.C) {C};

                  % Emitter Loop (Left) connected to Emitter (Left node)
                  % Horizontal branch L->R (to fix M1 orientation)
                  \draw (-6,0) to[switch, l=$S_1$] (-4.5,0)
                  to[R, l=$R_1$] (-3,0)
                  to[ammeter, l=$M_1$] (-1.5,0)
                  -- (Q1.E);

                  % Vertical branch T->B (Battery) and Ground return
                  % VEE: + to Emitter. Standard battery (+ Top).
                  \draw (-6,0) to[battery1, l=$V_{EE}$] (-6,-2.5)
                  -| (0,-2.5) % Connect to ground line
                  -- (Q1.B);

                  % Collector Loop (Right) connected to Collector (Right node)
                  % C is Negative relative to B.
                  % Path: C -> M2 -> R2 -> Ground.
                  % No VCC here (Fig 11-14).
                  \draw (Q1.C) -- (1.5,0)
                  to[ammeter, l=$M_2$] (3,0)
                  to[vR, l=$R_2$] (4.5,0)
                  |- (0,-2.5);

                  % Ground symbol at base connection (Bottom)
                  \draw (0,-2.5) node[ground] {};
              \end{circuitikz}
              \caption{Mediciones de corriente y tensión en el circuito emisor-base.}
              \label{fig:11-14}
          \end{figure}

    \item Aplicar la potencia. Observar y medir la corriente en el circuito de emisor y colector. Anotar
          los datos en la tabla~\ref{tab:11-1}. Medir con un EVM la tensión emisor-base y la de colector-base.
          Anotar los datos en la tabla~\ref{tab:11-1}. Indicar la polaridad de la tensión.
    \item Desconectar la alimentación de potencia. Modificar el circuito de la fig.~\ref{fig:11-14}
          añadiendo \( V_{CC} \) y \( S_2 \), como en la fig.~\ref{fig:11-15}. Esta segunda batería polariza
          el circuito colector-base.

          \begin{figure}[H]
              \centering
              \begin{circuitikz}[american, scale=1.0]
                  % PNP Transistor
                  \draw (0,0) node[pnp, rotate=90] (Q1) {};
                  \node[right] at (Q1.B) {B};
                  \node[below] at (Q1.E) {E};
                  \node[below] at (Q1.C) {C};

                  % Emitter Loop (Left)
                  \draw (-6,0) to[switch, l=$S_1$] (-4.5,0)
                  to[R, l=$R_1$] (-3,0)
                  to[ammeter, l=$M_1$] (-1.5,0)
                  -- (Q1.E);

                  \draw (-6,0) to[battery1, l=$V_{EE}$] (-6,-2.5) % Standard (+ top)
                  -| (0,-2.5)
                  -- (Q1.B);

                  % Collector Loop (Right) - Adds VCC and S2
                  % VCC: Negative to Collector.
                  % Path: C -> ... -> VCC -> Ground.
                  % Top of VCC connected to C. So Top must be Negative.
                  % Standard battery is + Top. So use 'invert'.
                  \draw (Q1.C) -- (1.5,0)
                  to[ammeter, l=$M_2$] (3,0)
                  to[vR, l=$R_2$] (4.5,0)
                  to[switch, l=$S_2$] (6,0)
                  to[battery1, l=$V_{CC}$, invert] (6,-2.5)
                  -| (0,-2.5);

                  \draw (0,-2.5) node[ground] {};
              \end{circuitikz}
              \caption{Circuito modificado con \( V_{CC} \) y \( S_2 \).}
              \label{fig:11-15}
          \end{figure}

    \item Cerrar \( S_2 \). Ajustar \( R_2 \) para la máxima resistencia. Cerrar \( S_1 \). Observar y
          medir las corrientes de emisor y de colector. Anotar los datos en la tabla~\ref{tab:11-1}. Medir y
          anotar en la tabla~\ref{tab:11-1} la tensión emisor-base y colector-base. Indicar la polaridad.
    \item Ajustar \( R_2 \) para la mínima resistencia. Conmutar los alcances del medidor como se
          requiera. Observar y medir la corriente de emisor y la de colector y las tensiones emisor-base y
          colector-base. Anotar los datos en la tabla~\ref{tab:11-1}. Indicar la polaridad de la tensión.
    \item Desconectar la alimentación. Invertir la polaridad de la batería de emisor (\( V_{EE} \)) y
          las conexiones del medidor \( M_1 \) (fig. 11-16). Ajustar \( R_2 \) para la mínima resistencia.

          \begin{figure}[H]
              \centering
              \begin{circuitikz}[american, scale=1.0]
                  % PNP Transistor
                  \draw (0,0) node[pnp, rotate=90] (Q1) {};
                  \node[right] at (Q1.B) {B};
                  \node[below] at (Q1.E) {E};
                  \node[below] at (Q1.C) {C};

                  % Emitter Loop (Left) - INVERTED VEE
                  % Inverted of Normal. Normal PNP VEE was + Top.
                  % So we want - Top. Use 'invert'.
                  \draw (-6,0) to[switch, l=$S_1$] (-4.5,0)
                  to[R, l=$R_1$] (-3,0)
                  to[ammeter, l=$M_1$] (-1.5,0)
                  -- (Q1.E);

                  \draw (-6,0) to[battery1, l=$V_{EE}$, invert] (-6,-2.5)
                  -| (0,-2.5)
                  -- (Q1.B);

                  % Collector Loop (Right) - Same as 11-15
                  \draw (Q1.C) -- (1.5,0)
                  to[ammeter, l=$M_2$] (3,0)
                  to[vR, l=$R_2$] (4.5,0)
                  to[switch, l=$S_2$] (6,0)
                  to[battery1, l=$V_{CC}$, invert] (6,-2.5)
                  -| (0,-2.5);

                  \draw (0,-2.5) node[ground] {};
              \end{circuitikz}
              \caption{Inversión de polaridad de \( V_{EE} \) y conexiones de \( M_1 \).}
              \label{fig:11-16}
          \end{figure}

    \item Aplicar la potencia. Observar y medir la corriente en los circuitos de emisor y de
          colector. Anotar los datos en la tabla~\ref{tab:11-1}. Medir y anotar en la tabla~\ref{tab:11-1}
          la tensión emisor-base y la tensión colector-base. Indicar la polaridad de la batería.
    \item Abrir \( S_2 \). Abrir el circuito emisor-base abriendo el interruptor \( S_1 \).
    \item Cerrar \( S_2 \). Observar y medir las corrientes de emisor y de colector. Anotar los datos
          en la tabla~\ref{tab:11-1}. Medir y anotar en la tabla~\ref{tab:11-1} la tensión colector-base.
          Desconectar la alimentación de potencia.
\end{enumerate}

\subsection{Polarización NPN}
\begin{enumerate}
    \setcounter{enumi}{10}
    \item Desconectar la alimentación. Quitar el transistor PNP del circuito y
          sustituirlo por uno NPN, como en la fig.~\ref{fig:11-14}. Invertir la polaridad de
          \( V_{CC} \) como en la fig.~\ref{fig:11-14}.

          \begin{figure}[H]
              \centering
              \begin{circuitikz}[american, scale=1.0]
                  % NPN Transistor
                  \draw (0,0) node[npn, rotate=90, yscale=-1] (Q1) {};
                  \node[right] at (Q1.B) {B};
                  \node[above] at (Q1.E) {E};
                  \node[above] at (Q1.C) {C};

                  % Emitter Loop (Left) - VEE Negative at Top for NPN Normal Bias
                  \draw (-6,0) to[switch, l=$S_1$] (-4.5,0)
                  to[R, l=$R_1$] (-3,0)
                  to[ammeter, l=$M_1$] (-1.5,0)
                  -- (Q1.E);

                  % Battery: invert for - Top
                  \draw (-6,0) to[battery1, l=$V_{EE}$, invert] (-6,-2.5)
                  -| (0,-2.5)
                  -- (Q1.B);

                  % Collector Loop (Right) - VCC Positive at Top for NPN Normal Bias
                  \draw (Q1.C) -- (1.5,0)
                  to[ammeter, l=$M_2$] (3,0)
                  to[vR, l=$R_2$] (4.5,0)
                  to[switch, l=$S_2$] (6,0)
                  % Battery: Normal for + Top
                  to[battery1, l=$V_{CC}$] (6,-2.5)
                  -| (0,-2.5);

                  \draw (0,-2.5) node[ground] {};
              \end{circuitikz}
              \caption{Circuito NPN con polarización normal.}
              \label{fig:11-17}
          \end{figure}
    \item Cerrar \( S_2 \). Ajustar \( R_2 \) para la máxima resistencia. Cerrar \( S_1 \). Observar
          y medir las corrientes de emisor y de colector. Anotar los datos en la tabla~\ref{tab:11-1}. Medir
          y anotar en la tabla~\ref{tab:11-1} la tensión emisor-base y colector-base. Indicar la polaridad.
    \item Ajustar \( R_2 \) para la mínima resistencia. Conmutar los alcances del medidor como se requiera.
          Observar y medir la corriente de emisor y la de colector y las tensiones emisor-base y colector-base.
          Anotar los datos en la tabla~\ref{tab:11-1}. Indicar la polaridad de la tensión.
    \item Desconectar la alimentación. Invertir la polaridad de la batería de emisor (\( V_{EE} \)) y las
          conexiones del medidor \( M_1 \) (fig. 11-16). Ajustar \( R_2 \) para la mínima resistencia.

          \begin{figure}[H]
              \centering
              \begin{circuitikz}[american, scale=1.0]
                  % NPN Transistor
                  \draw (0,0) node[npn, rotate=90, yscale=-1] (Q1) {};
                  \node[right] at (Q1.B) {B};
                  \node[above] at (Q1.E) {E};
                  \node[above] at (Q1.C) {C};

                  % Emitter Loop (Left) - VEE Positive at Top (Inverted Bias for NPN)
                  \draw (-6,0) to[switch, l=$S_1$] (-4.5,0)
                  to[R, l=$R_1$] (-3,0)
                  to[ammeter, l=$M_1$] (-1.5,0)
                  -- (Q1.E);

                  % Battery: Normal for + Top
                  \draw (-6,0) to[battery1, l=$V_{EE}$] (-6,-2.5)
                  -| (0,-2.5)
                  -- (Q1.B);

                  % Collector Loop (Right) - Same as 11-17
                  \draw (Q1.C) -- (1.5,0)
                  to[ammeter, l=$M_2$] (3,0)
                  to[vR, l=$R_2$] (4.5,0)
                  to[switch, l=$S_2$] (6,0)
                  to[battery1, l=$V_{CC}$] (6,-2.5)
                  -| (0,-2.5);

                  \draw (0,-2.5) node[ground] {};
              \end{circuitikz}
              \caption{Circuito NPN con polarización inversa de emisor.}
              \label{fig:11-18}
          \end{figure}

    \item Aplicar la potencia. Observar y medir la corriente en los circuitos de emisor y de colector.
          Anotar los datos en la tabla~\ref{tab:11-1}. Medir y anotar en la tabla~\ref{tab:11-1} la tensión
          emisor-base y la tensión colector-base. Indicar la polaridad de la batería.
    \item Abrir \( S_2 \). Abrir el circuito emisor-base abriendo el interruptor \( S_1 \).
    \item Cerrar \( S_2 \). Observar y medir las corrientes de emisor y de colector. Anotar los datos
          en la tabla~\ref{tab:11-1}. Medir y anotar en la tabla~\ref{tab:11-1} la tensión colector-base.
          Desconectar la alimentación de potencia.
\end{enumerate}

\subsection{Adicional}
\begin{enumerate}
    \setcounter{enumi}{17}
    \item La relación entre la corriente de colector y la corriente de emisor es:
          \begin{equation}
              I_C = I_E - I_B
          \end{equation}
    \item Verificar experimentalmente esta ecuación para transistores PNP y NPN. Dibujar el esquema del
          circuito utilizado y explicar detalladamente el procedimiento seguido. Consignar los resultados de
          las mediciones en una tabla especialmente preparada para ello.
\end{enumerate}

\begin{table}[H]
    \centering
    \resizebox{\textwidth}{!}{
        \begin{tabular}{|l|c|c|c|c|c|}
            \hline
            \textbf{Paso / Condición}               & \textbf{$I_E$ (mA)} & \textbf{$I_C$ (mA)} & \textbf{$V_{EB}$ (V)} & \textbf{$V_{CB}$ (V)} & \textbf{Polaridad} \\
            \hline
            1. PNP - Fig 11-14 (S1 cerrado)         &                     &                     &                       &                       &                    \\
            \hline
            2. PNP - Fig 11-15 (S2 cerrado, R2 máx) &                     &                     &                       &                       &                    \\
            \hline
            3. PNP - Fig 11-15 (S2 cerrado, R2 mín) &                     &                     &                       &                       &                    \\
            \hline
            4. PNP - Fig 11-16 (VEE invertida)      &                     &                     &                       &                       &                    \\
            \hline
            5. PNP - Fig 11-16 (S1 abierto)         &                     &                     &                       &                       &                    \\
            \hline
            6. NPN - Fig 11-17 (S2 cerrado, R2 máx) &                     &                     &                       &                       &                    \\
            \hline
            7. NPN - Fig 11-17 (S2 cerrado, R2 mín) &                     &                     &                       &                       &                    \\
            \hline
            8. NPN - Fig 11-16 mod (VEE invertida)  &                     &                     &                       &                       &                    \\
            \hline
            9. NPN - S1 abierto                     &                     &                     &                       &                       &                    \\
            \hline
        \end{tabular}
    }
    \caption{Tabla de datos}
    \label{tab:11-1}
\end{table}

\section{\textbf{Preguntas}}
\begin{enumerate}
    \item ¿Cual de las fig.~\ref{fig:11-14} a \ref{fig:11-18} muestra la polarización correcta de los circuitos de emisor y
          colector para:
          \begin{itemize}
              \item un transistor PNP.
              \item un transistor NPN.
          \end{itemize}
          Explicarlo refiriendose a los resultados consignados en la tabla~\ref{tab:11-1}.
    \item ¿Cual es el efecto sobre la corriente de colector de aumentar la polarización de emisor?
    \item ¿Cuales son los efectos sobre la corriente de colector de invertir la polarización en
          el circuito emisor-base? Refierase a sus propios datos para concretar su respuesta.
    \item ¿Que es Icbo? ¿Explique como la ha medido?
    \item ¿Que son los huecos? ¿En que tipo de material semiconductor los huecos son los portadores
          mayoritarios de corriente? ¿Por qué?
    \item ¿Que precauciones se deben observar cuando se efectún mediciones en los circuitos de
          transistores? Enuméralas.
\end{enumerate}
\begin{comment}
\section{Respuestas a las preguntas de autoexamen.}
\begin{enumerate}
    \item Silicio.
    \item Dos.
    \item En sentido directo; en sentido inverso.
    \item Positivo; negativo.
    \item Invertidas.
    \item Verdad.
    \item Falso.
    \item Verdad.
    \item Al calor; radiadores de calor.
    \item Tener cuidado; el deterioro.
\end{enumerate}
\end{comment}
\end{document}
